% Part: first-order-logic
% Chapter: models-theories
% Section: introduction

\documentclass[../../../include/open-logic-section]{subfiles}

\begin{document}

\olfileid{fol}{mat}{int}
\olsection{Inleiding}

\begin{explain}
De ontwikkeling van de axiomatische methode is een belangrijke
prestatie in de geschiedenis van de wetenschap en is van bijzonder
belang in de geschiedenis van de wiskunde. Bij de axiomatische opbouw
van een vakgebied moeten tal van vragen worden opgehelderd: Waarover
gaat het vakgebied? Wat zijn de meest fundamentele begrippen? Hoe
hangen ze samen? Kunnen alle begrippen van het vakgebied in termen van
deze fundamentele begrippen worden gedefinieerd? Welke wetten gelden
ervoor, en aan welke wetten moeten deze begrippen voldoen?

De axiomatische methode en de logica zijn voor elkaar gemaakt. De
formele logica biedt de middelen om axiomatische theorieën te
formuleren, om stellingen op nauwkeurig vastgelegde wijze uit de
axioma's van de theorie te bewijzen en om de eigenschappen van alle
systemen die aan de axioma's voldoen systematisch te bestuderen.
\end{explain}

\begin{defn}
Een verzameling !!{sentence}s~$\Gamma$ is \emph{gesloten} dan en
slechts dan als uit $\Gamma \Entails !A$ steeds $!A \in \Gamma$
volgt. De \emph{afsluiting} van een verzameling !!{sentence}s~$\Gamma$
is $\Setabs{!A}{\Gamma \Entails !A}$.

We zeggen dat~$\Gamma$ \emph{wordt geaxiomatiseerd door} een
verzameling zinnen~$\Delta$ als $\Gamma$ de afsluiting van~$\Delta$ is.
\end{defn}

\begin{explain}
Een axiomatische theorie kunnen we opvatten als de verzameling zinnen
die door haar verzameling axioma's~$\Delta$ wordt geaxiomatiseerd. Met
andere woorden: neem een eerste-ordetaal met niet-logische symbolen
voor de primitieve begrippen van de axiomatisch ontwikkelde wetenschap
die we willen bestuderen, en daarnaast een verzameling !!{sentence}s
die de fundamentele wetten van die wetenschap uitdrukken. Dan kunnen
we de theorie weergeven door alle !!{sentence}s in deze taal die
semantisch uit de axioma's volgen. Dit loopt uiteen van eenvoudige
voorbeelden met slechts één primitief begrip en eenvoudige axioma's,
zoals de theorie van partiële ordeningen, tot complexe theorieën zoals
de Newtoniaanse mechanica.

De volgende belangrijke logische feiten maken deze formele benadering
van de axiomatische methode zo belangrijk. Stel dat $\Gamma$ een
axiomasysteem voor een theorie is, dus een verzameling zinnen.
\begin{enumerate}
\item We kunnen precies aangeven wanneer een axiomasysteem een beoogde
  klasse !!{structure}s karakteriseert. Als we geïnteresseerd zijn in
  een bepaalde klasse !!{structure}s, karakteriseert een
  axiomasysteem~$\Gamma$ die klasse precies dan en slechts dan als de
  !!{structure}s in die klasse precies de~$\Struct M$ zijn waarvoor
  $\Sat{M}{\Gamma}$.
\item Dit kan mislukken doordat er een~$\Struct M$ bestaat waarvoor
  $\Sat{M}{\Gamma}$, terwijl~$\Struct M$ niet een van de beoogde
  !!{structure}s is. Dit kan aanleiding geven om axioma's toe te
  voegen die in~$\Struct M$ niet waar zijn.
\item Als we er ten minste in slagen dat~$\Gamma$ in alle beoogde
  !!{structure}s waar is, dan is een zin~$!A$ in alle beoogde
  !!{structure}s waar wanneer $\Gamma \Entails !A$. We kunnen dus
  logische middelen (zoals methoden voor formele !!{derivation})
  gebruiken om aan te tonen dat zinnen in alle beoogde !!{structure}s
  waar zijn: we hoeven slechts aan te tonen dat ze semantisch uit de
  axioma's volgen.
\item Soms hebben we geen beoogde !!{structure}s voor ogen, maar
  vertrekken we van de axioma's zelf. We beginnen dan met enkele
  primitieve begrippen die aan bepaalde wetten moeten voldoen, en
  leggen die wetten vast in een axiomasysteem. We willen meteen nagaan
  of de axioma's elkaar niet tegenspreken. Als ze dat wel doen, bestaat
  er geen interpretatie van de begrippen die aan al die wetten voldoet
  en hebben we een incoherente theorie opgezet. Dat dit niet gebeurt,
  kunnen we verifiëren door een model van~$\Gamma$ te vinden. Als onze
  theorie modellen heeft, kunnen we die bovendien met logische methoden
  onderzoeken en construeren.
\item Ook de onafhankelijkheid van de axioma's is een belangrijke
  vraag. Het kan gebeuren dat een van de axioma's eigenlijk een gevolg
  van de andere is en dus redundant is. We kunnen bewijzen dat een
  axioma~$!A$ in~$\Gamma$ redundant is door
  $\Gamma \setminus \{!A\} \Entails !A$ te bewijzen. We kunnen ook
  bewijzen dat een axioma niet redundant is door aan te tonen dat
  $(\Gamma \setminus \{!A\}) \cup \{\lnot !A\}$ vervulbaar is. Zo
  werd bijvoorbeeld aangetoond dat het parallellenaxioma onafhankelijk
  is van de overige axioma's van de meetkunde.
\item Een andere belangrijke vraag betreft de definieerbaarheid van
  begrippen in een theorie. De keuze van de taal bepaalt waaruit de
  modellen van een theorie bestaan. Niet elk aspect van een theorie
  hoeft echter afzonderlijk in haar modellen te worden weergegeven.
  Iedere ordening~$\le$ bepaalt bijvoorbeeld een bijbehorende strikte
  ordening~$<$---wanneer de ene is gegeven, kunnen we de andere
  definiëren. Een model van een theorie met zo'n ordening hoeft de
  bijbehorende strikte ordening dus niet \emph{ook} als afzonderlijke
  relatie te bevatten. Wanneer kan een relatie in het algemeen in
  termen van andere relaties worden gedefinieerd? Wanneer is dat
  onmogelijk en moeten we haar dus aan de primitieve begrippen van de
  taal toevoegen?
\end{enumerate}
\end{explain}


\end{document}
